% SOURCE_STATEMENT: "How can a hierarchical material be designed systematically to achieve a particular set of properties? How can different levels of structural hierarchy be bridged? What criteria should determine the characteristic length scales at all hierarchical levels? What is the theoretical basis for bottom-up materials design?"
% Preparation and final self-audit: CONTRIBUTING.md prompts C and D; source/audit artifacts retained in the local physics_statements_2026_10_11 batch.
% !TeX program = lualatex
\documentclass[11pt]{article}
\usepackage[margin=25mm]{geometry}
\usepackage{fontspec}
\setmainfont{Latin Modern Roman}
\usepackage{amsmath,amssymb,mathtools}
\usepackage{unicode-math}
\setmathfont{Latin Modern Math}
\usepackage{xurl,hyperref}
\hypersetup{colorlinks=true,urlcolor=blue}
\setcounter{secnumdepth}{0}
\setlength{\parindent}{0pt}
\setlength{\parskip}{5pt}
\setlength{\emergencystretch}{3em}
\newcommand{\sref}[2]{\hyperlink{source:#1:#2}{[S#2]}}
\newcommand{\cataloguescope}[1]{\paragraph{Recorded target.}#1}
\newcommand{\R}{\mathbb R}
\newcommand{\C}{\mathbb C}
\newcommand{\N}{\mathbb N}
\newcommand{\Z}{\mathbb Z}
\newcommand{\Q}{\mathbb Q}
\newcommand{\F}{\mathbb F}
\newcommand{\eps}{\varepsilon}
\begin{document}
\section{Systematic Multiscale Design of Hierarchical Materials}\label{statement}

\subsection{Definitions and mathematical statement}

\paragraph{Shared notation.}$\mathbb N=\{1,2,\ldots\}$; the standard number fields have their usual meanings. All additional conventions are stated locally.


\paragraph{Definitions and conventions.} A hierarchical material has structural features at several characteristic length scales $\ell_1<\cdots<\ell_m$, rather than a single microstructure scale. Effective properties, such as stiffness, strength, and flaw tolerance, are macroscopic responses derived from the constitutive and geometric behavior at finer scales. In the report, flaw tolerance means that a pre-existing crack does not propagate before the surrounding material approaches its limiting strength.

\paragraph{Source scope and qualifications.} This is a four-part design research program. The source does not specify a universal class of admissible microstructures, a property map, an optimization norm, or an algorithmic complexity target. These must be chosen for any subsequent precise model; adding them here would invent a narrower problem.

\paragraph{Catalogue mathematical target.}
\paragraph{Catalogue wording retained for traceability.}
How can a hierarchical material be designed systematically to achieve a particular set of properties? How can different levels of structural hierarchy be bridged? What criteria should determine the characteristic length scales at all hierarchical levels? What is the theoretical basis for bottom-up materials design?

\paragraph{Saved literature context.} The supplied catalogue's assessment dated 2026-08-21 records: Hierarchical and flaw-tolerant material models provide design principles and examples, but the general systematic multiscale design, scale-selection criterion, and theoretical bottom-up basis asked here remain open research questions. The assessment is retained as dated provenance; it is not a new comprehensive literature or open-status certification.
\subsection{Short English statement}
Find a systematic theory for designing material structure across many scales to achieve specified mechanical properties.
\subsection{Sources}
\begin{itemize}
\item[\hypertarget{source:statement:1}{S1}] ULAM.AI and the UnsolvedMath Contributors, UnsolvedMath, catalogue entry. Catalogue curation under CC BY 4.0; definitions and English exposition added locally.
\url{https://huggingface.co/datasets/ulamai/UnsolvedMath}.
\item[\hypertarget{source:statement:2}{S2}] Mini-Workshop: Mathematical Methods and Models of Continuum Biomechanics (2005), Oberwolfach Reports, DOI 10.4171/owr/2005/09, pp. 503--528; official report source and mathematical context.
\url{https://doi.org/10.4171/owr/2005/09}.
\item[\hypertarget{source:statement:3}{S3}] Official EMS Press PDF of the same Oberwolfach report; PDF page 6, printed page 508.
\url{https://ems.press/content/serial-article-files/45983}.
\item[\hypertarget{source:statement:4}{S4}] Gao et al., materials become insensitive to flaws at nanoscale; reference in the saved dated catalogue literature review.
\url{https://doi.org/10.1073/pnas.0937635100}.
\end{itemize}
\label{end:statement}
\end{document}
